\section{多模态：从视频理解到生成模型融合}

本章进入第四条主线：多模态。多模态的发展不是简单“给语言模型加图片”，而是视觉理解、视频建模、生成模型、视觉 Transformer、图文对齐和扩散 Transformer 逐步汇合。这里的关键词是融合：不同模态最终要在统一表征、统一生成或统一推理框架里互相影响。

\begin{figure}[H]
\centering
\includegraphics[width=0.96\textwidth]{multimodal-lineage.png}
\caption{多模态模型谱系：视频、GAN、Diffusion、ViT、CLIP、Stable Diffusion 与 DiT 汇合。自制概念图，依据 03:08:08--03:56:38 对谈内容整理。}
\end{figure}

\begin{knowledgebox}{读图：多模态不是一条直线}
视频理解、图像生成、视觉 Transformer、图文对齐和扩散模型分别从不同问题出发。它们后来汇合，是因为模型需要同时理解世界、生成世界，并用语言把视觉能力接入用户意图。
\end{knowledgebox}

\subsection{DeepVideo、双流网络与早期视频理解}

DeepVideo 和双流网络代表早期深度学习进入视频领域的尝试。图像可以被 CNN 处理之后，研究者自然会问：视频是不是也可以？难点在于视频多了时间维度，模型不仅要识别单帧内容，还要理解动作、运动和事件。

\begin{importantbox}{视频比图像多出的难题}
视频建模至少多了三个问题：时间依赖、运动表示和长程事件。单帧图像回答“这是什么”，视频还要回答“它怎样变化”“为什么这样变化”“接下来会怎样”。
\end{importantbox}

\subsection{GAN、Diffusion 与 DDPM}

图像生成线先由 GAN 主导。GAN 能生成锐利样本，但训练不稳定；VAE 训练稳定但图像容易模糊。Diffusion 早期在 GAN 阴影下成长，DDPM 让扩散模型重新回到图像舞台中央。扩散模型的优势在于训练稳定、生成质量好，并能和条件控制、文本引导结合。

\begin{knowledgebox}{GAN 与 Diffusion 的取舍}
\begin{tabularx}{\textwidth}{p{0.22\textwidth}p{0.34\textwidth}X}
\toprule
模型族 & 优势 & 难点 \\
\midrule
GAN & 样本锐利，生成速度快 & 训练不稳定，模式崩塌风险高。 \\
VAE & 训练稳定，有概率建模解释 & 样本容易模糊。 \\
Diffusion & 稳定、高质量、易条件控制 & 采样成本较高，需要工程优化。 \\
\bottomrule
\end{tabularx}
\end{knowledgebox}

\subsection{ViT、CLIP、Stable Diffusion 与 DiT}

前面讲生成模型如何走向扩散，本节看视觉和语言如何汇合到统一架构。ViT 把图像切成 patch，并用 Transformer 处理视觉 token；CLIP 把图像和文本放入同一语义空间，成为文生图和多模态检索的重要基石；Stable Diffusion 把扩散模型、潜空间和开源生态结合起来，推动图像生成普及；DiT 则把 Diffusion 与 Transformer 更深地融合，代表人们对统一架构未来的期待。

\begin{importantbox}{CLIP 为什么关键}
CLIP 的重要性在于图文对齐。它让模型知道一张图和一段文本在语义上是否匹配，从而支撑文生图、图像检索、数据过滤和多模态表示学习。没有图文对齐，生成模型很难稳定听懂用户提示。
\end{importantbox}

\subsection{本章小结}

多模态线说明，视觉、语言和生成正在靠拢。GAN 解决生成起点，Diffusion 解决质量和稳定，ViT 把视觉拉进 Transformer，CLIP 建立图文桥梁，DiT 则把生成和 Transformer 进一步合流。

